\phantomsection\addcontentsline{toc}{section}{Drones}
\ECchapter{18}{Drones}{How do autonomous aerial systems remain stable and safe?}{ECOrange}

\ECObjectives{
\begin{itemize}
\item Explain how thrust, attitude and position are controlled in a multirotor drone.
\item Relate sensor characteristics and control-loop bandwidth to flight stability.
\item Build a safe mission plan from payload, endurance, weather and fault-response requirements.
\end{itemize}}

\begin{multicols}{2}
\begin{engineeringnews}
Drones are increasingly used for inspection, mapping, emergency response and precision agriculture.
Their usefulness depends on more than flight performance: safe autonomy requires reliable sensing,
energy management, communication and clearly defined procedures for fault conditions.
\end{engineeringnews}

\begin{ecarticle}
A multirotor drone flies by controlling the thrust produced by several propellers. When the total
upward thrust equals the weight, the aircraft hovers. Increasing all rotor speeds causes a vertical
acceleration. Changing the balance between rotors creates roll, pitch or yaw moments and allows the
drone to move horizontally or rotate.

The flight controller closes several feedback loops. Gyroscopes measure angular velocity and
accelerometers measure specific acceleration. Their signals are combined to estimate attitude. A
barometer provides altitude information, while satellite navigation and cameras can estimate
position and horizontal velocity. Sensor fusion is essential because every sensor has limitations:
gyroscopes drift, accelerometers are disturbed by vibration and GNSS signals may be inaccurate or
unavailable indoors.

A common controller compares the desired attitude with the estimated attitude and adjusts motor
speeds many times per second. The inner loop stabilises angular motion, while slower outer loops
control velocity and position. If tuning is poor, the aircraft may oscillate or react too slowly to a
gust. Sampling frequency, actuator response and filtering must therefore be considered together.

Endurance is strongly limited by battery mass. A larger battery stores more energy but also increases
weight and the power required to hover. Payload has the same effect. Engineers therefore optimise the
frame, propellers, motors and flight plan together. Wind, temperature and repeated acceleration also
reduce real endurance compared with ideal calculations.

Autonomy introduces additional requirements. The drone must detect obstacles, maintain a safe
separation from people and execute a predictable action if communication is lost. Depending on the
mission, this may be hovering, landing or returning to a predefined point. Geofencing can prevent
entry into restricted zones, but it does not replace pilot responsibility, maintenance or risk
assessment.
\end{ecarticle}

\begin{tcolorbox}[title=\sffamily\bfseries Engineering Insight,colback=yellow!10,colframe=ECOrange]
A quadrotor normally changes direction by tilting. Part of the total thrust then acts horizontally,
so the motors must generate additional thrust to maintain altitude. Aggressive manoeuvres therefore
consume energy faster than steady flight even when the travelled distance is unchanged.
\end{tcolorbox}

\begin{engineeringfocus}
\textbf{Nested control loops of a drone}
\begin{center}
\resizebox{.98\linewidth}{!}{%
\begin{tikzpicture}[font=\scriptsize,>=Latex,node distance=3.5mm]
\node[draw=ECBlue,fill=ECLightBlue,rounded corners,align=center] (cmd) {position\\command};
\node[draw=ECGreen,fill=ECGreen!10,rounded corners,align=center,right=of cmd] (pos) {position\\controller};
\node[draw=ECOrange,fill=ECOrange!10,rounded corners,align=center,right=of pos] (att) {attitude\\controller};
\node[draw=ECPurple,fill=ECPurple!10,rounded corners,align=center,below=of att] (mot) {motors and\\propellers};
\node[draw=ECRed,fill=ECRed!8,rounded corners,align=center,below=of pos] (sens) {IMU, barometer,\\GNSS, camera};
\draw[->,thick] (cmd)--(pos); \draw[->,thick] (pos)--(att); \draw[->,thick] (att)--(mot);
\draw[->,thick] (mot)--(sens); \draw[->,thick] (sens)--(pos); \draw[->,thick] (sens)--(att);
\end{tikzpicture}%
}
\end{center}
Fast attitude control stabilises the aircraft; slower position control generates the attitude command.
\end{engineeringfocus}

\begin{keyfigures}
\begin{tabularx}{\linewidth}{@{}Xr@{}}
Hover condition & thrust = weight\\
Fast sensor & gyroscope\\
Altitude sensor & barometer\\
Main energy store & battery\\
Typical fault action & return or land
\end{tabularx}
\end{keyfigures}

\begin{tcolorbox}[title=\sffamily\bfseries Payload and endurance,colback=white,colframe=ECOrange]
\begin{center}
\begin{tikzpicture}
\begin{axis}[width=.94\linewidth,height=49mm,xlabel={Payload (kg)},ylabel={Endurance (min)},xmin=0,xmax=1.2,
ymin=10,ymax=36,grid=major,ticklabel style={font=\scriptsize}]
\addplot[very thick,mark=*] table[x=payload,y=endurance,col sep=comma]{data/EC18/drone_endurance.csv};
\end{axis}
\end{tikzpicture}
\end{center}
\footnotesize Illustrative trend for one aircraft: endurance decreases non-linearly as payload
increases. The graph must not be extrapolated beyond the validated operating range.
\end{tcolorbox}

\begin{vocabulary}
\begin{tabularx}{\linewidth}{@{}lX@{}}
thrust & poussée\\
hover & vol stationnaire\\
roll / pitch / yaw & roulis / tangage / lacet\\
flight controller & contrôleur de vol\\
gyroscope & gyroscope\\
sensor fusion & fusion de capteurs\\
payload & charge utile\\
gust & rafale\\
geofencing & géorepérage\\
return-to-home & retour automatique
\end{tabularx}
\end{vocabulary}

\begin{applicationexercise}
The endurance graph gives 25 min for a payload of 0.6 kg. A safety rule requires the drone to land
with 20\% of the initial flight time still available. The outward journey to an inspection area takes
6 min and the return journey takes the same time.
\begin{enumerate}
\item Calculate the maximum authorised flight time before landing.
\item Determine the time available for inspection over the target area.
\item Recalculate this inspection time for a payload of 0.9 kg.
\item Apply an additional 15\% endurance reduction for cold weather and recalculate the mission margin.
\item State whether the mission should be accepted and justify your decision.
\end{enumerate}
\end{applicationexercise}

\begin{questions}
\item How does a quadrotor produce roll and pitch motion?
\item Why is sensor fusion necessary?
\item Distinguish the attitude loop from the position loop.
\item Explain the compromise created by a larger battery.
\item What should happen after a communication failure?
\item Why should a control loop not use an excessively noisy measurement directly?
\end{questions}

\begin{engineeringchallenge}
Define a drone mission to inspect a bridge. Write requirements for payload, image resolution,
endurance, wind and separation distance. Select the main sensors, propose a safe path and define the
response to low battery, GNSS loss and communication loss. Include an energy budget, a pre-flight
checklist and three validation tests before operational approval.
\end{engineeringchallenge}

\begin{speaking}
Should drones be authorised to deliver parcels in cities? Compare efficiency, noise, safety, privacy
and infrastructure requirements, then state one quantitative acceptance criterion.
\end{speaking}
\end{multicols}

\subsection*{Sources}
\ECsource{European Union Aviation Safety Agency -- Drones}{https://www.easa.europa.eu/en/domains/civil-drones-rpas}
\ECsource{NASA -- Unmanned Aircraft Systems}{https://www.nasa.gov/centers-and-facilities/ames/unmanned-aircraft-systems/}